% =====================================================================
%  Chapter: Derivative-Constrained Deep Smoothing
% =====================================================================
\chapter{Derivative-Constrained Deep Smoothing}
\label{ch:deepsmoother}

The parametric benchmark of \cref{ch:parametric} quantified a trade-off:
per-slice SVI fits well but leaks static arbitrage; SSVI is arbitrage-free by
construction but three times less accurate. This chapter implements a neural
smoother designed to occupy the space between, combining the
\emph{architecture} of \cite{ackerer2020deep} with the \emph{training
mechanics} of \cite{hoshisashi2024noarbitrage}.

\section{Method}
\label{sec:dsmethod}

\subsection{Architecture: an SSVI prior times a neural corrector}
Following \cite{ackerer2020deep}, the total-variance surface is modelled
multiplicatively,
\begin{equation}
  w_\theta(k,\tau) \;=\; w_{\mathrm{SSVI}}(k,\tau)\,\times\,
  \mathcal C_\theta(k,\tau),
  \qquad
  \mathcal C_\theta = \exp\!\big(\mathrm{MLP}_\theta(k,\tau)\big) > 0,
  \label{eq:priorcorrector}
\end{equation}
where the prior is an SSVI surface with a smooth power-law term structure
$\theta(\tau)=\alpha\tau^{\beta}$ fitted to the day's quotes, and the
corrector is a small multilayer perceptron ($2\times64$ hidden units,
$\tanh$ activations) whose inputs $(k,\tau)$ are rescaled to $[-1,1]^2$.
With standard small-weight initialization $\mathcal C_\theta\approx1$, so
training \emph{starts at the prior} and departs from it only where the data
demand. The choice of $\tanh$ is not cosmetic: the loss below involves
$\partial^2_k w_\theta$, and piecewise-linear activations such as ReLU have
vanishing second derivatives almost everywhere, which would silence the
butterfly penalty \cite{hoshisashi2024noarbitrage}.

\subsection{Loss: soft no-arbitrage penalties on a collocation grid}
Training minimizes
\begin{equation}
  \mathcal L(\theta)
  = \underbrace{\frac1N\sum_{i=1}^{N}\big(w_\theta(k_i,\tau_i)-w_i\big)^2}_{\text{fit on sparse quotes}}
  + \lambda_{\mathrm{bfly}}\,\overline{\relu\!\big(-g_\theta\big)^2}\Big|_{\hat X}
  + \lambda_{\mathrm{cal}}\,\overline{\relu\!\big(-\partial_\tau w_\theta\big)^2}\Big|_{\hat X},
  \label{eq:dsloss}
\end{equation}
where $g_\theta$ is the Durrleman function of \cref{eq:durrleman} evaluated
on the model surface, and --- the decisive detail imported from
\cite{hoshisashi2024noarbitrage} --- the penalties are averaged over a
\emph{dense collocation grid} $\hat X$ ($24\times10$ points spanning the
day's quoted rectangle) that is disjoint from the quotes. No-arbitrage is
thus enforced everywhere, including regions with no data, which is what makes
extrapolation safe. The derivatives $w_k,w_{kk},w_\tau$ entering
\cref{eq:dsloss} are computed by \emph{exact} automatic differentiation of
the network (nested \texttt{elementwise\_grad}); before any training we
verified the machinery against the closed-form SVI derivatives of
\cref{ch:foundations}, with agreement at the $10^{-16}$ level. Soft
penalties, rather than hard architectural constraints, follow the evidence of
\cite{chataigner2020deep} that hard constraints sacrifice too much
representational power; we set $\lambda_{\mathrm{bfly}}=\lambda_{\mathrm{cal}}=10$.

\section{Controlled experiments on synthetic surfaces}
\label{sec:dsablation}

The method is first validated where ground truth is known: a target SSVI
surface sampled on sparse irregular grids ($6$--$14$ quotes per slice, $1.5\%$
noise), with a deliberately misspecified prior fitted from the same sparse
quotes. Three ablations isolate the two design choices.

\begin{table}[t]
  \centering
  \caption[Deep smoother ablations]{Ablation study on a synthetic surface
  (IV RMSE against the dense ground truth, in vol points; $\min g$ and
  violation share on the collocation grid).}
  \label{tab:ablation}
  \begin{tabular}{lS[table-format=1.3]S[table-format=+1.3]S[table-format=2.1]}
    \toprule
    Variant & {RMSE (vol pts)} & {$\min g$} & {Violations (\%)} \\
    \midrule
    Full model (prior + constraints) & 0.187 & +0.33 & 0.0 \\
    No constraints ($\lambda=0$)     & 0.247 & {---} & {---} \\
    No prior (flat BS prior)         & 1.418 & {---} & {---} \\
    \bottomrule
  \end{tabular}
  \\[2pt]\footnotesize \TODO{fill the $\min g$/violation columns of the
  ablated variants from the notebook run.}
\end{table}

Two lessons emerge from \cref{tab:ablation}. First, the \emph{prior is the
dominant ingredient} in the sparse-data regime: without it the corrector must
invent the entire surface shape and the error grows sevenfold. Second, the
constraints act as a \emph{regularizer} as well as a safety device --- the
fully constrained model fits \emph{better} than the unconstrained one, in
line with the observations of \cite{hoshisashi2024noarbitrage}, while
guaranteeing a clean $g$ on the collocation grid
(\cref{fig:gheatmap}).

\begin{figure}[t]
  \centering
  \maybegraphics[width=\textwidth]{figures/nb03_gheatmap.png}
  \caption[Arbitrage heat-map of the deep smoother]{Durrleman $g(k,\tau)$
  heat-map. Left: the constrained model is non-negative everywhere. Right:
  removing the penalties lets butterfly arbitrage (red) leak into regions
  unsupported by data.}
  \label{fig:gheatmap}
\end{figure}

\begin{figure}[t]
  \centering
  \maybegraphics[width=0.8\textwidth]{figures/nb03_tradeoff.png}
  \caption[Training dynamics: fit vs penalty]{Training dynamics of the fit
  term and butterfly penalty (log scale), in the style of Fig.~7 of
  \cite{hoshisashi2024noarbitrage}.}
  \label{fig:tradeoff}
\end{figure}

\paragraph{Robustness to sparsity.} Subsampling the synthetic quotes exposes
the structural advantage of a \emph{surface-level} smoother over per-slice
calibration: at $25\%$ of the quotes, no slice retains the six points needed
to calibrate a five-parameter SVI --- per-slice calibration simply ceases to
exist --- while the deep smoother, which pools information across maturities,
still reconstructs the surface at $0.315$ vol points
(\cref{fig:sparsity}).

\begin{figure}[t]
  \centering
  \maybegraphics[width=0.7\textwidth]{figures/nb03_sparsity.png}
  \caption[Robustness to quote sparsity]{Reconstruction error as quotes are
  subsampled. Per-slice SVI loses calibrable slices and eventually cannot be
  evaluated; the surface-level smoother degrades gracefully.}
  \label{fig:sparsity}
\end{figure}

\section{Real-data protocol and preliminary results}
\label{sec:dsreal}

On real SPX days the protocol mirrors the benchmark exactly: OTM quotes
grouped by expiration, an SSVI prior fitted per day, $20\%$ of quotes
withheld, penalties on the day's collocation grid, and the same
arbitrage audit. On the preliminary sample the smoother attains a hold-out
error of $0.158$ vol points against $0.153$ in-sample --- hold-out
$\approx$ in-sample, as for the parametric fits, with $\min g>0$ on the
audit grid. \TODO{replace with full-sample per-day results joined against
\texttt{benchmark\_svi\_slices\_\{full,paper\}.parquet}, and add the
comparison figure deep-vs-SVI-vs-SSVI per day.}

\begin{figure}[t]
  \centering
  \maybegraphics[width=\textwidth]{figures/nb03_smiles.png}
  \caption[Deep smoother on real SPX quotes]{The deep smoother on a real
  trading day: quotes (dots) and fitted smiles per maturity.}
  \label{fig:dssmiles}
\end{figure}

\section{Discussion}
The prior-times-corrector design earns both of its components: the prior
carries the sparse-data regime, the derivative penalties buy near-exact
no-arbitrage at zero cost in accuracy --- on synthetic data they even improve
it. The remaining honest caveat, inherited from the soft-constraint
paradigm and documented by \cite{chataigner2020deep}, is that penalties
\emph{reduce} violations without \emph{abolishing} them on unseen regions;
we therefore always report the residual violation rate alongside the error.
The chapter's limitation is economic rather than statistical: one network
must be trained \emph{per day}. Removing that cost is the subject of
\cref{ch:operator}.
